% vim: ai sw=2 spell spelllang=cs

\begin{frame}[fragile]
  \frametitle{ELF}

  \heading{Executable and Linkable Format}
  \begin{itemize}
    \item Hlavní souborový formát na Linuxu
    \item Přenositelný, relokovatelný, rozšiřitelný

    \pause

    \item Podpora pro ladicí informace
    \begin{itemize}
      \item \iffimuni Přespříště \else V další sekci \fi
    \end{itemize}

    \pause

    \item \cmd{readelf} je specializovaný \cmd{objdump} pro ELF
    \item \cmd{libelf} je specializovaná \cmd{libbfd} pro ELF

    \pause

    \item Literatura
    \begin{itemize}
      \item \href{http://www.sco.com/developers/devspecs/gabi41.pdf}{System V Application Binary Interface}
    \end{itemize}
  \end{itemize}

\end{frame}

\begin{frame}{\file{libelf}}{Srovnání}

  \begin{itemize}
    \item \file{libbfd} je nezávislá na binárním formátu
    \begin{itemize}
      \item Ale neumí vše, co podporuje ELF
    \end{itemize}

    \pause

    \item \file{libelf} umí jen ELF
    \begin{itemize}
      \item Ale zvládá téměř všechno, co ELF podporuje
      \item Pracuje na nižší úrovni
    \end{itemize}

    \pause

    \item \file{libebl} je nadstavba \file{libelf}
    \begin{itemize}
      \item Abstraktnější rozhraní
      \item Nemá moc uživatelů
    \end{itemize}
  \end{itemize}

\end{frame}

\begin{frame}{\file{libelf}}{Detaily}

  \begin{itemize}
    \item Knihovna pro práci s ELF
    \item Nabízí dvě API
    \begin{itemize}
      \item ELF (\header{libelf.h})
      \item GELF -- více abstraktní, generické (\header{gelf.h})
      \item Lze použít současně a míchat
    \end{itemize}

    \pause

    \item Dokumentace
    \begin{itemize}
      \item Pro BSD port: \href{http://sourceforge.net/projects/elftoolchain/files/Documentation/libelf-by-example/20120308/libelf-by-example.pdf/download}{A tutorial introduction to libelf}
      \item V hlavičkových souborech
    \end{itemize}

    \pause

    \item Knihovny při překladu: \cmd{gcc ... -lelf}
    \item První je třeba volat \code{elf_version(EV_CURRENT)}
    \begin{itemize}
      \item Nesmí vrátit \code{EV_NONE}
    \end{itemize}
  \end{itemize}

\end{frame}

\begin{frame}[fragile]
  \frametitle{\file{libelf}}
  \framesubtitle{Ošetření chyb}

  \heading{Jak zjistit, co se stalo?}
  \begin{itemize}
    \item Poslední chyba: \code{int elf_errno()}
    \item Převod na text: \code{const char *elf_errmsg(int error)}
    \begin{itemize}
      \item \code{\-1} jako \code{error} je zkratka pro \code{elf_errno()}
    \end{itemize}
  \end{itemize}

  \begin{block}{Typicky}
  \begin{lstlisting}
if (!elf_function(...))
  errx(1, "elf_function: %s", elf_errmsg(-1));
  \end{lstlisting} %stopzone
  \end{block}

\end{frame}

\begin{frame}
  \frametitle{Úkol}

  \heading{Inicializace \file{libelf} a otevření souboru}
  \begin{enumerate}
    \item Najděte a otevřete si BSD dokumentaci
    \item Upravujte \file{pb173/04/}
    \item Zavolejte \code{elf_version(EV_CURRENT)}
    \item Zkontrolujte návratovou hodnotu
    \begin{itemize}
      \item Nesmí být \code{EV_NONE}
      \item V případě chyby chybu vypište (\code{elf_errmsg(-1)})
    \end{itemize}

    \item Pomocí standardního \code{open} otevřete soubor z příkazové řádky
    \begin{itemize}
      \item \code{argv[1]}
      \item Pro čtení
    \end{itemize}

    \item Přeložte a spusťte
  \end{enumerate}

\end{frame}

\begin{frame}[fragile]{\file{libelf}}{Soubory}

  \begin{itemize}
    \item \file{libelf} pracuje s \emph{file deskriptory}

    \pause

    \item Začátek práce
    \begin{itemize}
      \item \code{Elf *elf_begin(int fd, Elf_Cmd cmd, Elf *ref)}
      \item \code{cmd} je jedno z \code{ELF_C_READ}, \code{ELF_C_WRITE},
	\code{ELF_C_RDWR}
      \item \code{ref} je \code{NULL}
    \end{itemize}

    \pause

    \item Zjištění typu sobouru
    \begin{itemize}
      \item \code{Elf_Kind elf_kind(Elf *elf)}
      \item Návratová hodnota je jedno z \code{ELF_K_NONE}, \code{ELF_K_AR},
              \emph{\code{ELF_K_ELF}}
    \end{itemize}

    \pause

    \item Konec práce
    \begin{itemize}
      \item \code{int elf_end(Elf *elf)}
    \end{itemize}
  \end{itemize}

  \pause

  \begin{block}{Příklad (bez ověření chyb)}
  \begin{lstlisting}[aboveskip=2pt,belowskip=2pt]
int fd = open(file, O_RDONLY);
Elf *elf = elf_begin(fd, ELF_C_READ, NULL);
elf_end(elf);
close(fd);
  \end{lstlisting}
  \end{block}

\end{frame}

\begin{frame}{Úkol}

  \heading{Doplňte načtení souboru zadaného jako parametr}
  \begin{enumerate}
    \item Zavolejte \code{elf_begin}

    \item Zavolejte \code{elf_kind}
    \begin{itemize}
      \item Ověřte, že soubor je typu \code{ELF_K_ELF}
      \item Pokud ne, program ukončete s chybou
    \end{itemize}

    \item Zavolejte \code{elf_end}
    \item Ověřujte návratové hodnoty

    \item Přeložte a spusťte
    \begin{itemize}
      \item S nějakým ELFem
      \item S nějakou statickou knihovnou
    \end{itemize}
  \end{enumerate}

\end{frame}

% vim: ai sw=2 spell spelllang=cs

\begin{frame}{Struktura ELFu}

  \begin{itemize}
    \item ELF hlavička

    \pause

    \item Hlavičky sekcí
    \begin{itemize}
      \item Popis dat pro linkování, ladění apod.
      \item Pro překladač, debugger, \dots
    \end{itemize}

    \item Hlavičky programové
    \begin{itemize}
      \item Pohled na data pro spuštění
      \item Pro interpretr
    \end{itemize}

    \pause

    \item Data
    \begin{itemize}
      \item Odkazovány z obou typů hlaviček
    \end{itemize}
  \end{itemize}

  \begin{textblock*}{5cm}(\textwidth-6cm,.5\textheight-2.3cm)
    % vim: ai sw=2 spell spelllang=cs

\tikzstyle{pre} = [<-,>=stealth',semithick]
\tikzstyle{sec} = [draw,semithick,minimum width=3.5cm, minimum height=8mm]
\begin{tikzpicture}[node distance=1pt, font=\scriptsize]
  \node [sec] (HEAD) {ELF header};
  \node [sec,below=of HEAD] (PRGHDR) {Program Header Table};
  \node [sec,below=of PRGHDR] (TEXT) {.text};
  \node [sec,below=of TEXT] (RODATA) {.rodata};
  \node [minimum width=3.5cm,below=of RODATA] (MORE) {\dots};
  \node [sec,below=of MORE] (DATA) {.data};
  \node [sec,below=of DATA] (SECHDR) {Section Header Table}
    edge [pre,->,in=350,out=4,looseness=2.2] node {} (DATA)
    edge [pre,->,in=350,out=0] node {} (RODATA)
    edge [pre,->,in=350,out=356] node {} (TEXT);

  \draw [semithick,decorate,decoration={brace,amplitude=6pt}]
    (RODATA.south west) -- (TEXT.north west)
    node [midway,xshift=-3pt,yshift=-1pt] {}
    edge [pre,in=182,out=160,looseness=1.5] node {} (PRGHDR);

  \draw [semithick,decorate,decoration={brace,amplitude=6pt}]
    (DATA.south west) -- (MORE.north west)
    node [midway,xshift=-3pt,yshift=-1pt] {}
    edge [pre,in=178,out=160,looseness=1.2] node {} (PRGHDR);
\end{tikzpicture}

  \end{textblock*}

\end{frame}



\begin{frame}
  \frametitle{ELF hlavička}

  \begin{itemize}
    \item \cmd{readelf -h}

    \pause

    \item Magické číslo (\code{0x7f 'E' 'L' 'F'})

    \pause

    \item Typ, cílová architektura, stroj, systém, \dots

    \pause

    \item Metadata popisující hlavičky

    \pause

    \item V \file{libelf} (\header{gelf.h})
    \begin{itemize}
      \item Struktura \code{GElf_Ehdr}
      \item \code{GElf_Ehdr *gelf_getehdr(Elf *elf, GElf_Ehdr *dest)}
    \end{itemize}
  \end{itemize}

\end{frame}

\begin{frame}
  \frametitle{Úkol}

  \heading{Výpis ELF hlavičky}
  \begin{enumerate}
    \item Vypište si ELF hlavičku pomocí \cmd{readelf}
    \item Ze svého programu vypište
    \begin{itemize}
      \item Velikost ELF hlavičky
      \item Architekturu, pro kterou ELF je
      \item Počet programových hlaviček
      \item Počet hlaviček sekcí
    \end{itemize}

    \item Přeložte a spusťte
  \end{enumerate}

\end{frame}

\begin{frame}
  \frametitle{Sekce}

  \begin{itemize}
    \item \cmd{readelf -S} (obsah sekce: \cmd{readelf -x},
      popř.\ text \cmd{readelf -p})

    \pause

    \item Vytvářené překladačem/linkerem
    \item Čtené překladačem/linkerem/interpretrem

    \pause

    \item Předdefinované sekce:
    \begin{itemize}
      \item \code{.text}: kód
      \item \code{.data}: nekonstantní data (proměnné)
      \item \code{.rodata}: konstantní data (řetězce apod.)
      \item \code{.bss}: data, která se inicializují při startu na 0
      \item \code{.interp}: interpretr
      \item \code{.symtab}: tabulka symbolů pro ladění (\cmd{strip})
      \item \code{.dynsym}: tabulka symbolů pro dynamický linker
      \item \code{.gnu_debuglink}: odkaz na soubor s ladicími informacemi
    \end{itemize}
  \end{itemize}

\end{frame}

\begin{frame}[fragile]
  \frametitle{Úkol}

  \heading{Práce se sekcemi}
  \begin{enumerate}
    \item Vytvořte si soubor s kódem \code{puts("hello")} a \code{puts("world")}
    \item Vytvořte si vlastní sekci s daty
    \begin{lstlisting}[belowskip=0pt]
int __attribute__((section(".data.my_section"))) x = 3;
    \end{lstlisting}
    \item Vytvořte si vlastní sekci s funkcí (\code{.text.my_section})
    \item Přeložte do \file{.o} a poté i slinkujte
    \begin{itemize}
      \item $\Rightarrow$ mějte 2 soubory
    \end{itemize}

    \item Vypište si seznam sekcí v obou souborech (\cmd{readelf -S})
    \item Vypište si obsah sekcí (\cmd{readelf -x}):
    \begin{itemize}
      \item \code{.rodata} (zkuste i \cmd{readelf -p})
      \item V \file{.o}: \code{.data.my_section}
      \item V \file{.o}: \code{.text.my_section} (zkuste i \cmd{objdump -d})
    \end{itemize}
  \end{enumerate}

  \medskip

  Pozn.: tyto soubory nezahazujte.

\end{frame}

\begin{frame}{\file{libelf}}{Sekce}

  \begin{itemize}
    \item Držátko: struktura \code{Elf_Scn}
    \begin{itemize}
      \item Obsah nedefinovaný
    \end{itemize}

    \pause

    \item Hlavička: struktura \code{GElf_Shdr} (\header{gelf.h})
    \begin{itemize}
      \item Získání hlavičky:
	\code{GElf_Shdr *gelf_getshdr(Elf_Scn *scn, GElf_Shdr *dst)}

      \pause

      \item \code{sh_name}: offset do dat sekce \code{.shstrtab}
      \item \code{sh_type}: \code{SHT_NULL}, \code{SHT_PROGBITS},
	\code{SHT_SYMTAB}, \code{SHT_STRTAB}, \dots
      \item \code{sh_flags}: \code{SHF_WRITE}, \code{SHF_ALLOC},
	\code{SHF_EXECINSTR}, \dots
      \item \code{sh_size}: velikost
    \end{itemize}

    \pause

    \item Data: struktura \code{Elf_Data} (obsah probereme \iffimuni příště\else
      dále\fi)
    \begin{itemize}
      \item Iterace přes obsah:
	\code{Elf_Data *elf_getdata(Elf_Scn *scn, Elf_Data *data)}

      \pause

      \item \code{d_buf}: data
      \item \code{d_type}: \code{ELF_T_BYTE}, \code{ELF_T_ADDR}, \dots
      \item \code{d_size}: velikost dat
    \end{itemize}
  \end{itemize}

\end{frame}

\begin{frame}[fragile]{\file{libelf}}{Vyhledání sekcí}

  \begin{itemize}
    \item x-tá sekce: \code{Elf_Scn *elf_getscn(Elf *elf, size_t index)}
    \item Iterace: \code{Elf_Scn *elf_nextscn(Elf *elf, Elf_Scn *scn)}

    \pause

    \item Index sekce s názvy sekcí \code{.shstrtab}: \\
      \code{int elf_getshdrstrndx(Elf *elf, size_t *dst)}
  \end{itemize}

  \onslide<2->

  \begin{block}{Příklad}
  \begin{lstlisting}[aboveskip=2pt,belowskip=2pt]
Elf_Scn *scn = NULL;
while ((scn = elf_nextscn(elf, scn))) {
  GElf_Shdr shdr;
  Elf_Data *data = NULL;

  gelf_getshdr(scn, &shdr);
  /* shdr */

  while ((data = elf_getdata(scn, data)))
    /* data */;
}
  \end{lstlisting}
  \end{block}

\end{frame}

\begin{frame}
  \frametitle{Úkol}

  \heading{Hexdump sekcí}
  \begin{enumerate}
    \item Iterujte přes sekce (\code{elf_nextscn})
    \item Vypište hlavičku každé sekce (\code{gelf_getshdr})
    \begin{itemize}
      \item Index (\code{elf_ndxscn()})
      \item Offset do sekce názvů (\code{sh_name})
      \item Velikost v hexa (\code{sh_}\dots)
      \item Flagy v hexa (\code{sh_}\dots)
    \end{itemize}

    \item Vypište data každé sekce (\code{elf_getdata()})
    \begin{itemize}
      \item Uvažte jen první \code{elf_getdata} (tj.\ neiterujte přes data)
      \item Hexdump prvních 16 bytů obsahu
    \end{itemize}

    \item Přeložte a vyzkoušejte
    \item Porovnejte s \cmd{readelf -S}
  \end{enumerate}

\end{frame}

\begin{frame}[fragile]{\file{libelf}}{Úpravy ELFu}

  \begin{itemize}
    \item Otevření s \code{ELF_C_WRITE}/\code{ELF_C_RDWR}

    \pause

    \item Nová sekce: \code{Elf_Scn *elf_newscn(Elf *elf)}
    \begin{itemize}
      \item \code{gelf_getshdr} jako předtím a naplnění
      \item Po naplnění: \code{int gelf_update_shdr(Elf_Scn *scn, GElf_Shdr *src)}
      \item Vytvoření dat: \code{Elf_Data *elf_newdata(Elf_Scn *scn)}
      \begin{itemize}
	\item I vícekrát -- sloučí se
      \end{itemize}
    \end{itemize}

    \pause

    \item Zápis do souboru
    \begin{itemize}
      \item \code{loff_t elf_update(Elf *elf, Elf_Cmd cmd)}
      \item \code{cmd}: \code{ELF_C_WRITE}
    \end{itemize}
  \end{itemize}

\end{frame}

\begin{frame}{Úkol}

  \heading{Přidání sekce do ELFu/ověření (ne)funkčnosti \file{libelf}}
  \begin{enumerate}
    \item Otevřete ELF soubor pro čtení a zápis
    \item Vytvořte novou sekci \code{.comment.my}
    \begin{itemize}
      \item Typu \code{SHT_NOTE}
    \end{itemize}
    \item Vložte do ní nějaká data
    \begin{itemize}
      \item Typu \code{ELF_T_BYTE}
    \end{itemize}
    \item Zavolejte \code{elf_update}
    \begin{itemize}
      \item S parametrem \code{ELF_C_WRITE}
    \end{itemize}
    \item Přeložte a vyzkoušejte
    \item Zobrazte obsah sekce \code{.comment.my} pomocí \cmd{readelf -x}
  \end{enumerate}

\end{frame}
